\section{The tangent space of the Picard functor: the truncated exponential}
\label{pa-chap:Tangent}

Kleiman's computation of the tangent space of the Picard scheme at the identity
\cite{kleiman-picard},
\(T_0 \Pic_{X/k} = H^1(X, \struct X)\), rests on the
\emph{truncated exponential sequence}: over the dual numbers
\(R[\varepsilon] = R \oplus R\varepsilon\) (\(\varepsilon^2 = 0\)), the units
\(1 + b\varepsilon\) form a copy of the additive group of \(R\) inside
\(R[\varepsilon]^\times\), split off by reduction mod \(\varepsilon\). This chapter
develops that layer in pure commutative algebra --- the unit group of the dual numbers,
its functoriality, the \(\varepsilon \mapsto a\varepsilon\) scaling, dual numbers under
base change, and the resulting two-chart \v{C}ech computation
\[
  H^1(X, \struct{X}) \;\cong\;
    \ker\Bigl( \check{H}^1{}^\times\bigl(\Gamma(U_0 \cap U_1)[\varepsilon]\bigr)
      \longrightarrow \check{H}^1{}^\times\bigl(\Gamma(U_0 \cap U_1)\bigr) \Bigr)
\]
for a scheme covered by two affine opens --- the two-cover \v{C}ech form of
\(H^1(C, \struct C) \cong \ker(\Pic(C_\varepsilon) \to \Pic(C))\).

(Here and below we write \(R[\varepsilon]\) for the dual numbers of a commutative ring
\(R\), and, for an element \(x = a + b\varepsilon\), \(\mathrm{fst}(x) = a\) and
\(\mathrm{snd}(x) = b\) for its two components.)

\section{Units of the dual numbers and the truncated exponential}

\begin{definition}[The truncated exponential unit]
  \label{pa-def:truncExpUnit}
  \lean{TruncExpCech.inr_mul_inr_eq_zero, TruncExpCech.truncExpUnit,
    TruncExpCech.truncExpUnit_val, TruncExpCech.fst_truncExpUnit,
    TruncExpCech.snd_truncExpUnit, TruncExpCech.truncExpUnit_zero}
  \leanok
  \source{kleiman-picard}

  \dcref{custom}
  Let \(R\) be a commutative ring and \(b \in R\). Since
  \((b\varepsilon)(c\varepsilon) = 0\) in \(R[\varepsilon]\), the element
  \(1 + b\varepsilon\) is a unit of \(R[\varepsilon]\), with inverse
  \(1 - b\varepsilon\). We write
  \(\mathrm{tExp}(b) := 1 + b\varepsilon \in R[\varepsilon]^\times\); its components are
  \(\mathrm{fst} = 1\) and \(\mathrm{snd} = b\), and \(\mathrm{tExp}(0) = 1\).
\end{definition}

\begin{lemma}[The truncated exponential is an injective homomorphism]
  \label{pa-lem:truncExp_add}
  \lean{TruncExpCech.truncExpUnit_add, TruncExpCech.truncExpUnit_injective,
    TruncExpCech.truncExp, TruncExpCech.truncExp_apply}
  \leanok
  \dcref{custom}

  \uses{pa-def:truncExpUnit}
  The map \(\mathrm{tExp} : (R, +) \to R[\varepsilon]^\times\) is a group homomorphism,
  \[
    \mathrm{tExp}(b + c) = \mathrm{tExp}(b) \cdot \mathrm{tExp}(c),
  \]
  and it is injective.
\end{lemma}
\begin{proof}
  \leanok
  \((1 + b\varepsilon)(1 + c\varepsilon) = 1 + (b + c)\varepsilon +
  bc\,\varepsilon^2 = 1 + (b+c)\varepsilon\). Injectivity: the second component of
  \(\mathrm{tExp}(b)\) is \(b\).
\end{proof}

\begin{definition}[Reduction mod \(\varepsilon\) on units, and its section]
  \label{pa-def:unitsFst}
  \lean{TruncExpCech.fstRingHom, TruncExpCech.fstRingHom_apply,
    TruncExpCech.inlRingHom, TruncExpCech.inlRingHom_apply,
    TruncExpCech.unitsFst, TruncExpCech.unitsFst_apply_val,
    TruncExpCech.unitsInl, TruncExpCech.unitsInl_apply_val}
  \leanok
  \dcref{custom}

  The component projection \(\mathrm{fst} : R[\varepsilon] \to R\),
  \(a + b\varepsilon \mapsto a\), and the constant inclusion
  \(\mathrm{inl} : R \to R[\varepsilon]\), \(a \mapsto a + 0\varepsilon\), are ring
  homomorphisms. They induce group homomorphisms on unit groups,
  \[
    \mathrm{fst}^\times : R[\varepsilon]^\times \to R^\times
    \qquad\text{and}\qquad
    \mathrm{inl}^\times : R^\times \to R[\varepsilon]^\times .
  \]
\end{definition}

\begin{lemma}[Splitting identities]
  \label{pa-lem:unitsFst_split}
  \lean{TruncExpCech.unitsFst_unitsInl, TruncExpCech.unitsFst_truncExpUnit,
    TruncExpCech.fst_inv_mul_fst, TruncExpCech.fst_mul_fst_inv}
  \leanok
  \dcref{custom}

  \uses{pa-def:unitsFst, pa-def:truncExpUnit}
  \(\mathrm{fst}^\times \circ \mathrm{inl}^\times = \mathrm{id}\) on \(R^\times\);
  \(\mathrm{fst}^\times(\mathrm{tExp}(b)) = 1\) for every \(b \in R\); and for any unit
  \(u \in R[\varepsilon]^\times\),
  \(\mathrm{fst}(u^{-1}) \cdot \mathrm{fst}(u) = 1 = \mathrm{fst}(u) \cdot
  \mathrm{fst}(u^{-1})\).
\end{lemma}
\begin{proof}
  \leanok
  The first two are componentwise computations
  (\(\mathrm{fst}(\mathrm{inl}(a)) = a\), \(\mathrm{fst}(1 + b\varepsilon) = 1\)); the
  third applies the ring homomorphism \(\mathrm{fst}\) to
  \(u^{-1} u = 1 = u\,u^{-1}\).
\end{proof}

\begin{theorem}[Exactness in the middle]
  \label{pa-thm:truncExp_exact}
  \lean{TruncExpCech.truncExp_range_eq_ker_unitsFst}
  \leanok
  \source{kleiman-picard}

  \uses{pa-lem:truncExp_add, pa-lem:unitsFst_split}
  \dcref{custom}
  The range of the truncated exponential is exactly the kernel of reduction mod
  \(\varepsilon\) on units: a unit \(u \in R[\varepsilon]^\times\) satisfies
  \(\mathrm{fst}^\times(u) = 1\) if and only if \(u = 1 + b\varepsilon\) for a (unique)
  \(b \in R\). Together with \ref{pa-lem:unitsFst_split} this exhibits the split short
  exact sequence
  \[
    1 \longrightarrow (R, +) \xrightarrow{\ \mathrm{tExp}\ } R[\varepsilon]^\times
      \xrightarrow{\ \mathrm{fst}^\times\ } R^\times \longrightarrow 1 .
  \]
\end{theorem}
\begin{proof}
  \leanok
  One inclusion is \ref{pa-lem:unitsFst_split}. Conversely, if
  \(\mathrm{fst}^\times(u) = 1\), then \(u = 1 + b\varepsilon\) for
  \(b = \mathrm{snd}(u)\), and \(u = \mathrm{tExp}(b)\); uniqueness is the injectivity
  of \ref{pa-lem:truncExp_add}.
\end{proof}

\begin{lemma}[The unit decomposition of the dual numbers]
  \label{pa-lem:unit_decomposition}
  \lean{TruncExpCech.unitsInl_unitsFst_mul_truncExpUnit}
  \leanok
  \dcref{custom}

  \uses{pa-def:unitsFst, pa-def:truncExpUnit, pa-lem:unitsFst_split}
  Every unit \(u \in R[\varepsilon]^\times\) factors as a constant unit times a
  truncated exponential:
  \[
    u \;=\; \mathrm{inl}^\times\bigl(\mathrm{fst}^\times(u)\bigr)
      \cdot \mathrm{tExp}\bigl(\mathrm{fst}(u^{-1}) \cdot \mathrm{snd}(u)\bigr).
  \]
\end{lemma}
\begin{proof}
  \leanok
  Compare components. First components: \(\mathrm{fst}(u) \cdot 1 = \mathrm{fst}(u)\).
  Second components: the product's second component is
  \(\mathrm{fst}(u) \cdot \mathrm{fst}(u^{-1}) \cdot \mathrm{snd}(u) =
  \mathrm{snd}(u)\) by \ref{pa-lem:unitsFst_split}.
\end{proof}

\begin{definition}[Functoriality of the dual numbers]
  \label{pa-def:dualNumber_map}
  \lean{TruncExpCech.mapRingHom, TruncExpCech.fst_mapRingHom,
    TruncExpCech.snd_mapRingHom}
  \leanok
  \dcref{custom}

  A ring homomorphism \(f : R \to S\) induces the ring homomorphism
  \[
    f[\varepsilon] : R[\varepsilon] \to S[\varepsilon], \qquad
    a + b\varepsilon \longmapsto f(a) + f(b)\varepsilon,
  \]
  applying \(f\) to both components.
\end{definition}

\begin{lemma}[Functoriality identities]
  \label{pa-lem:dualNumber_map_functorial}
  \lean{TruncExpCech.mapRingHom_id, TruncExpCech.mapRingHom_comp,
    TruncExpCech.mapRingHom_inl, TruncExpCech.unitsMap_mapRingHom_unitsInl}
  \leanok
  \dcref{custom}

  \uses{pa-def:dualNumber_map, pa-def:unitsFst}
  \(\mathrm{id}[\varepsilon] = \mathrm{id}\);
  \((g \circ f)[\varepsilon] = g[\varepsilon] \circ f[\varepsilon]\);
  \(f[\varepsilon] \circ \mathrm{inl} = \mathrm{inl} \circ f\); and on unit groups
  \(f[\varepsilon]^\times \circ \mathrm{inl}^\times = \mathrm{inl}^\times \circ
  f^\times\).
\end{lemma}
\begin{proof}
  \leanok
  All four are componentwise computations from Definition~\ref{pa-def:dualNumber_map}.
\end{proof}

\begin{lemma}[Naturality of the truncated exponential and of the reduction]
  \label{pa-lem:truncExp_naturality}
  \lean{TruncExpCech.map_mapRingHom_truncExpUnit, TruncExpCech.unitsFst_map_mapRingHom}
  \leanok
  \dcref{custom}

  \uses{pa-def:dualNumber_map, pa-def:truncExpUnit, pa-def:unitsFst}
  For a ring homomorphism \(f : R \to S\):
  \(f[\varepsilon]^\times(\mathrm{tExp}(b)) = \mathrm{tExp}(f(b))\) for all
  \(b \in R\), and
  \(\mathrm{fst}^\times \circ f[\varepsilon]^\times = f^\times \circ
  \mathrm{fst}^\times\) on \(R[\varepsilon]^\times\).
\end{lemma}
\begin{proof}
  \leanok
  Componentwise: \(f[\varepsilon](1 + b\varepsilon) = 1 + f(b)\varepsilon\), and
  \(\mathrm{fst}(f[\varepsilon](x)) = f(\mathrm{fst}(x))\).
\end{proof}

\begin{definition}[The \(\varepsilon \mapsto a\varepsilon\) scaling]
  \label{pa-def:mumford_scaling}
  \lean{TruncExpCech.scaleRingHom, TruncExpCech.fst_scaleRingHom,
    TruncExpCech.snd_scaleRingHom}
  \leanok
  \source{kleiman-picard}

  \dcref{custom}
  For \(a \in R\), the \emph{scaling} \(\mu_a : R[\varepsilon] \to R[\varepsilon]\) is
  the ring homomorphism \(r + m\varepsilon \mapsto r + (am)\varepsilon\) (the identity
  on constants, multiplication by \(a\) on the infinitesimal part). Scalar
  multiplication on the tangent space of a functor at a rational point is
  functoriality along \(\mu_a\).
\end{definition}

\begin{lemma}[Scaling identities]
  \label{pa-lem:mumford_scaling_props}
  \lean{TruncExpCech.scaleRingHom_comp_scaleRingHom, TruncExpCech.scaleRingHom_one,
    TruncExpCech.unitsScale_truncExpUnit, TruncExpCech.mapRingHom_comp_scaleRingHom}
  \leanok
  \dcref{custom}

  \uses{pa-def:mumford_scaling, pa-def:truncExpUnit, pa-def:dualNumber_map}
  \(\mu_a \circ \mu_b = \mu_{ab}\); \(\mu_1 = \mathrm{id}\); on truncated exponentials
  \(\mu_a^\times(\mathrm{tExp}(b)) = \mathrm{tExp}(ab)\); and for a ring homomorphism
  \(\rho : R \to S\) and \(s \in R\),
  \(\rho[\varepsilon] \circ \mu_s = \mu_{\rho(s)} \circ \rho[\varepsilon]\).
\end{lemma}
\begin{proof}
  \leanok
  All componentwise computations.
\end{proof}

\section{Dual numbers under base change}

\begin{definition}[The dual-number map over an algebra map]
  \label{pa-def:dualNumber_mapAlgHom}
  \lean{TruncExpCech.mapAlgHom, TruncExpCech.mapAlgHom_apply}
  \leanok
  \dcref{custom}

  \uses{pa-def:dualNumber_map}
  For a \(k\)-algebra \(A\), the functorial map
  \((\text{algebra map})[\varepsilon] : k[\varepsilon] \to A[\varepsilon]\) of
  Definition~\ref{pa-def:dualNumber_map} is a \(k\)-algebra homomorphism (the
  compatibility with the structure constants holds componentwise).
\end{definition}

\begin{theorem}[Dual numbers under base change]
  \label{pa-thm:dualNumber_baseChange}
  \lean{TruncExpCech.baseChangeAlgHom, TruncExpCech.baseChangeAlgHom_tmul,
    TruncExpCech.baseChangeAlgHom_bijective, TruncExpCech.baseChangeAlgEquiv,
    TruncExpCech.baseChangeAlgEquiv_tmul, TruncExpCech.baseChangeAlgEquiv_symm_apply}
  \leanok
  \dcref{custom}

  \uses{pa-def:dualNumber_mapAlgHom}
  Let \(k\) be a commutative ring and \(A\) a \(k\)-algebra. Extension of scalars of
  the dual numbers is the dual numbers of \(A\): there is an \(A\)-algebra isomorphism
  \[
    A \otimes_k k[\varepsilon] \;\cong\; A[\varepsilon], \qquad
    a \otimes (c + d\varepsilon) \longmapsto a \cdot (c + d\varepsilon)
      = ac + ad\,\varepsilon,
  \]
  with inverse \(x \mapsto \mathrm{fst}(x) \otimes 1 + \mathrm{snd}(x) \otimes
  \varepsilon\).
\end{theorem}
\begin{proof}
  \leanok
  The displayed map is the \(A\)-algebra extension of the \(k\)-algebra map of
  \ref{pa-def:dualNumber_mapAlgHom} along the scalar extension \(k \to A\); it is a
  well-defined \(A\)-algebra homomorphism. The candidate inverse
  \(g(x) = \mathrm{fst}(x) \otimes 1 + \mathrm{snd}(x) \otimes \varepsilon\) is
  additive. Composing one way,
  \(g(ac + ad\varepsilon) = ac \otimes 1 + ad \otimes \varepsilon =
  a \otimes (c + d\varepsilon)\) after moving scalars across the tensor
  (\(ac \otimes 1 = a \otimes c\cdot 1\) etc.), and every element of
  \(A \otimes_k k[\varepsilon]\) is a sum of pure tensors, on which this computation
  applies; composing the other way,
  \(\mathrm{fst}(x)\cdot 1 + \mathrm{snd}(x)\cdot\varepsilon = x\). So the map is
  bijective. (The dual numbers are finite free over \(k\), so no flatness input is
  needed --- the inverse is written down explicitly.)
\end{proof}

\section{The two-chart \v{C}ech unit-cocycle engine}

A \emph{two-chart datum} is a pair of ring homomorphisms \(\rho_1 : A_1 \to B\),
\(\rho_2 : A_2 \to B\) --- for the section rings of a two-affine cover of a scheme,
\(A_i = \Gamma(U_i)\), \(B = \Gamma(U_0 \cap U_1)\) and the \(\rho_i\) the restrictions.
Everything in this section is elementary commutative algebra over such a datum.

\begin{definition}[The multiplicative \v{C}ech coboundaries]
  \label{pa-def:cech_coboundary_units}
  \lean{TruncExpCech.cechCoboundaryUnits, TruncExpCech.mem_cechCoboundaryUnits,
    TruncExpCech.unitsMap_mem_cechCoboundaryUnits_left,
    TruncExpCech.unitsMap_mem_cechCoboundaryUnits_right}
  \leanok
  \dcref{custom}

  For a two-chart datum \(\rho_1 : A_1 \to B\), \(\rho_2 : A_2 \to B\), the
  \emph{coboundary subgroup} \(\mathcal B^\times(\rho_1, \rho_2) \leq B^\times\) is the
  subgroup generated by the images of \(\rho_1^\times\) and \(\rho_2^\times\); its
  elements are exactly the products \(\rho_1^\times(v_1)\,\rho_2^\times(v_2)\) with
  \(v_i \in A_i^\times\) (the two images are subgroups, so the products already form a
  subgroup). The quotient \(B^\times / \mathcal B^\times(\rho_1, \rho_2)\) is the
  two-cover \v{C}ech \(\check H^1\) of units --- for a two-affine cover, the \v{C}ech
  Picard group of the cover.
\end{definition}

\begin{definition}[The additive \v{C}ech coboundaries]
  \label{pa-def:cech_coboundary_add}
  \lean{TruncExpCech.cechCoboundaryAdd, TruncExpCech.mem_cechCoboundaryAdd,
    TruncExpCech.sub_mem_cechCoboundaryAdd,
    TruncExpCech.exists_sub_of_mem_cechCoboundaryAdd}
  \leanok
  \dcref{custom}

  For a two-chart datum as above, the \emph{additive coboundary subgroup}
  \(\mathcal B^+(\rho_1, \rho_2) \leq (B, +)\) is
  \(\rho_1(A_1) + \rho_2(A_2)\); its elements are exactly the sums
  \(\rho_1(a_1) + \rho_2(a_2)\), equivalently the differences
  \(\rho_1(a_1) - \rho_2(a_2)\) (replace \(a_2\) by \(-a_2\)). The quotient
  \(B / \mathcal B^+(\rho_1, \rho_2)\) is the two-cover \v{C}ech cokernel --- for the
  section rings of a two-affine cover, the carrier of \ref{pa-def:TwoCover_H1Cok}.
\end{definition}

\begin{theorem}[The truncated exponential detects the additive coboundaries]
  \label{pa-thm:truncExp_detects_coboundaries}
  \lean{TruncExpCech.truncExpUnit_mem_cechCoboundaryUnits_iff}
  \leanok
  \source{kleiman-picard}

  \uses{pa-def:cech_coboundary_units, pa-def:cech_coboundary_add, pa-lem:unit_decomposition,
    pa-lem:dualNumber_map_functorial, pa-lem:truncExp_naturality, pa-lem:truncExp_add,
    pa-lem:unitsFst_split}
  \dcref{custom}
  Let \(\rho_1, \rho_2\) be a two-chart datum and \(b \in B\). The truncated
  exponential \(\mathrm{tExp}(b) \in B[\varepsilon]^\times\) is a coboundary of the
  thickened datum \(\rho_1[\varepsilon], \rho_2[\varepsilon]\) if and only if \(b\) is
  an additive coboundary:
  \[
    \mathrm{tExp}(b) \in \mathcal B^\times\bigl(\rho_1[\varepsilon],
      \rho_2[\varepsilon]\bigr)
    \iff
    b \in \mathcal B^+(\rho_1, \rho_2).
  \]
\end{theorem}
\begin{proof}
  \leanok
  (\(\Leftarrow\)) If \(b = \rho_1(a_1) + \rho_2(a_2)\), then by naturality
  (\ref{pa-lem:truncExp_naturality}) and additivity (\ref{pa-lem:truncExp_add}),
  \(\mathrm{tExp}(b) = \mathrm{tExp}(\rho_1(a_1)) \cdot \mathrm{tExp}(\rho_2(a_2)) =
  \rho_1[\varepsilon]^\times(\mathrm{tExp}(a_1)) \cdot
  \rho_2[\varepsilon]^\times(\mathrm{tExp}(a_2))\), a coboundary.

  (\(\Rightarrow\)) Suppose \(\mathrm{tExp}(b) =
  \rho_1[\varepsilon]^\times(w_1) \cdot \rho_2[\varepsilon]^\times(w_2)\) with
  \(w_i \in A_i[\varepsilon]^\times\). Decompose each chart unit
  (\ref{pa-lem:unit_decomposition}) as
  \(w_i = \mathrm{inl}^\times(v_i)\cdot\mathrm{tExp}(c_i)\) with
  \(v_i \in A_i^\times\), \(c_i \in A_i\). Pushing through
  \(\rho_i[\varepsilon]\) (\ref{pa-lem:dualNumber_map_functorial},
  \ref{pa-lem:truncExp_naturality}) and collecting constant and exponential factors
  (truncated exponentials commute with everything and multiply additively),
  \[
    \mathrm{tExp}(b) = \mathrm{inl}^\times\bigl(\rho_1^\times(v_1)\rho_2^\times(v_2)\bigr)
      \cdot \mathrm{tExp}\bigl(\rho_1(c_1) + \rho_2(c_2)\bigr).
  \]
  Applying \(\mathrm{fst}^\times\) (\ref{pa-lem:unitsFst_split},
  \ref{pa-lem:truncExp_naturality}) kills both exponentials and gives
  \(1 = \rho_1^\times(v_1)\rho_2^\times(v_2)\); substituting back and using
  injectivity of \(\mathrm{tExp}\) (\ref{pa-lem:truncExp_add}) yields
  \(b = \rho_1(c_1) + \rho_2(c_2) \in \mathcal B^+(\rho_1, \rho_2)\).
\end{proof}

\begin{definition}[The \v{C}ech units reduction]
  \label{pa-def:cech_units_reduction}
  \lean{TruncExpCech.cechCoboundaryUnits_le_comap_unitsFst,
    TruncExpCech.cechUnitsReduction}
  \leanok
  \dcref{custom}

  \uses{pa-def:cech_coboundary_units, pa-lem:truncExp_naturality,
    pa-lem:dualNumber_map_functorial}
  Reduction mod \(\varepsilon\) on units carries thickened coboundaries to
  coboundaries: by naturality (\ref{pa-lem:truncExp_naturality},
  \ref{pa-lem:dualNumber_map_functorial}),
  \(\mathrm{fst}^\times\bigl(\rho_1[\varepsilon]^\times(w_1)\,
  \rho_2[\varepsilon]^\times(w_2)\bigr) =
  \rho_1^\times(\mathrm{fst}^\times w_1)\, \rho_2^\times(\mathrm{fst}^\times w_2)\).
  Hence \(\mathrm{fst}^\times\) descends to a group homomorphism of \v{C}ech
  \(\check H^1\)-of-units groups,
  \[
    \pi \;:\;
    B[\varepsilon]^\times \big/ \mathcal B^\times(\rho_1[\varepsilon],
      \rho_2[\varepsilon])
    \longrightarrow
    B^\times \big/ \mathcal B^\times(\rho_1, \rho_2),
  \]
  the \v{C}ech-cocycle incarnation of the restriction
  \(\Pic(X \times_k \Spec k[\varepsilon]) \to \Pic(X)\) along
  \(\varepsilon \mapsto 0\).
\end{definition}

\begin{lemma}[Truncated-exponential classes lie in the kernel, and vanish exactly on
  coboundaries]
  \label{pa-lem:cech_kernel_truncExp_class}
  \lean{TruncExpCech.cechUnitsReduction_mk_truncExpUnit,
    TruncExpCech.mk_truncExpUnit_eq_one_iff}
  \leanok
  \dcref{custom}

  \uses{pa-def:cech_units_reduction, pa-thm:truncExp_detects_coboundaries,
    pa-lem:unitsFst_split}
  For \(b \in B\), the class \([\mathrm{tExp}(b)]\) lies in \(\ker \pi\); and
  \([\mathrm{tExp}(b)] = 1\) in the thickened \v{C}ech group if and only if
  \(b \in \mathcal B^+(\rho_1, \rho_2)\).
\end{lemma}
\begin{proof}
  \leanok
  \(\pi[\mathrm{tExp}(b)] = [\mathrm{fst}^\times(\mathrm{tExp}(b))] = [1]\)
  (\ref{pa-lem:unitsFst_split}). Vanishing of the class means membership of
  \(\mathrm{tExp}(b)\) in the thickened coboundary subgroup, which is
  \ref{pa-thm:truncExp_detects_coboundaries}.
\end{proof}

\begin{theorem}[Every kernel class is a truncated exponential]
  \label{pa-thm:cech_kernel_surjectivity}
  \lean{TruncExpCech.exists_mk_truncExpUnit_of_cechUnitsReduction_eq_one}
  \leanok
  \source{kleiman-picard}

  \uses{pa-def:cech_units_reduction, pa-thm:truncExp_exact, pa-lem:dualNumber_map_functorial,
    pa-def:cech_coboundary_units}
  \dcref{custom}
  Every class \(x\) of
  \(B[\varepsilon]^\times / \mathcal B^\times(\rho_1[\varepsilon],
  \rho_2[\varepsilon])\) with \(\pi(x) = 1\) is represented by a truncated exponential:
  \(x = [\mathrm{tExp}(b)]\) for some \(b \in B\).
\end{theorem}
\begin{proof}
  \leanok
  Let \(u \in B[\varepsilon]^\times\) represent \(x\). Since \(\pi(x) = 1\), the
  reduction \(\mathrm{fst}^\times(u)\) is a coboundary:
  \(\mathrm{fst}^\times(u) = \rho_1^\times(v_1)\,\rho_2^\times(v_2)\) with
  \(v_i \in A_i^\times\). Correct \(u\) by the constant lifts of these chart units:
  \[
    u' := \Bigl(\rho_1[\varepsilon]^\times\bigl(\mathrm{inl}^\times v_1\bigr)\Bigr)^{-1}
      \cdot u \cdot
      \Bigl(\rho_2[\varepsilon]^\times\bigl(\mathrm{inl}^\times v_2\bigr)\Bigr)^{-1}.
  \]
  The two correction factors are thickened coboundaries
  (\ref{pa-def:cech_coboundary_units}), so \([u'] = [u] = x\); and by
  \ref{pa-lem:dualNumber_map_functorial} the reduction of \(u'\) is
  \(\rho_1^\times(v_1)^{-1} \cdot \mathrm{fst}^\times(u) \cdot \rho_2^\times(v_2)^{-1}
  = 1\). By exactness (\ref{pa-thm:truncExp_exact}), \(u' = \mathrm{tExp}(b)\) for some
  \(b \in B\).
\end{proof}

\begin{theorem}[The truncated-exponential kernel computation]
  \label{pa-thm:truncExp_cech_kernel_equiv}
  \lean{TruncExpCech.truncExpCechKernelMonoidHom,
    TruncExpCech.truncExpCechKernelAddEquiv,
    TruncExpCech.truncExpCechKernelAddEquiv_apply_mk}
  \leanok
  \source{kleiman-picard}

  \uses{pa-lem:cech_kernel_truncExp_class, pa-thm:cech_kernel_surjectivity, pa-lem:truncExp_add,
    pa-def:cech_coboundary_add}
  \dcref{custom}
  For a two-chart datum \(\rho_1 : A_1 \to B\), \(\rho_2 : A_2 \to B\), the truncated
  exponential \(b \mapsto [\mathrm{tExp}(b)]\) induces an isomorphism of abelian groups
  \[
    B \big/ \bigl(\rho_1(A_1) + \rho_2(A_2)\bigr)
    \;\cong\;
    \ker\Bigl( \pi :
      B[\varepsilon]^\times \big/ \mathcal B^\times(\rho_1[\varepsilon],
        \rho_2[\varepsilon])
      \to B^\times \big/ \mathcal B^\times(\rho_1, \rho_2) \Bigr),
  \]
  sending the residue class of \(b\) to the class of \(1 + b\varepsilon\).
\end{theorem}
\begin{proof}
  \leanok
  The assignment \(b \mapsto [\mathrm{tExp}(b)]\) is a homomorphism from \((B, +)\) into
  \(\ker \pi\) (\ref{pa-lem:truncExp_add} for additivity,
  \ref{pa-lem:cech_kernel_truncExp_class} for landing in the kernel). It kills exactly
  \(\mathcal B^+(\rho_1, \rho_2)\) (\ref{pa-lem:cech_kernel_truncExp_class}), so it
  descends to an injective homomorphism on the quotient
  \(B/(\rho_1(A_1) + \rho_2(A_2))\). It is surjective onto \(\ker \pi\) by
  \ref{pa-thm:cech_kernel_surjectivity}.
\end{proof}

\begin{lemma}[Scaling equivariance of the engine]
  \label{pa-lem:cech_scaling_equivariance}
  \lean{TruncExpCech.cechCoboundaryUnits_le_comap_unitsScale,
    TruncExpCech.unitsScale_mk_truncExpUnit}
  \leanok
  \dcref{custom}

  \uses{pa-def:mumford_scaling, pa-lem:mumford_scaling_props, pa-def:cech_coboundary_units,
    pa-def:truncExpUnit}
  Let \(\rho_1, \rho_2\) be a two-chart datum and let \(t \in B\) be
  \emph{chart-compatible}: \(t = \rho_1(s_1) = \rho_2(s_2)\) for some \(s_i \in A_i\).
  Then the scaling \(\mu_t^\times\) of \(B[\varepsilon]^\times\) preserves the
  thickened coboundary subgroup \(\mathcal B^\times(\rho_1[\varepsilon],
  \rho_2[\varepsilon])\), hence descends to the thickened \v{C}ech group; and on
  truncated-exponential classes it acts by
  \([\mathrm{tExp}(b)] \mapsto [\mathrm{tExp}(t \cdot b)]\).
\end{lemma}
\begin{proof}
  \leanok
  For a coboundary \(\rho_1[\varepsilon]^\times(w_1)\,\rho_2[\varepsilon]^\times(w_2)\),
  the compatibility \(\rho_i[\varepsilon] \circ \mu_{s_i} = \mu_{\rho_i(s_i)} \circ
  \rho_i[\varepsilon] = \mu_t \circ \rho_i[\varepsilon]\)
  (\ref{pa-lem:mumford_scaling_props}) rewrites its image under \(\mu_t^\times\) as
  \(\rho_1[\varepsilon]^\times(\mu_{s_1}^\times w_1)\,
  \rho_2[\varepsilon]^\times(\mu_{s_2}^\times w_2)\), again a coboundary. The action on
  truncated exponentials is \(\mu_t^\times(\mathrm{tExp}(b)) = \mathrm{tExp}(t b)\)
  (\ref{pa-lem:mumford_scaling_props}), passed to classes.
\end{proof}

\section{\texorpdfstring{\(H^1\)}{H1} as the
  \texorpdfstring{\(\varepsilon\)}{epsilon}-kernel on a two-cover}

We now assemble the engine on the pinned two-cover section rings of a scheme \(X\) over
\(\Spec k\) with two affine opens \(U_0 \cup U_1 = X\): the two-chart datum
is \(B = \Gamma(X, U_0 \cap U_1)\) with the two restrictions
\(\rho_i : \Gamma(X, U_i) \to B\).

\begin{lemma}[Module quotients are additive-subgroup quotients]
  \label{pa-lem:submodule_quotient_addEquiv}
  \lean{TruncExpCech.submoduleQuotientAddEquiv,
    TruncExpCech.submoduleQuotientAddEquiv_mk}
  \leanok
  \dcref{custom}

  For a module \(M\) over a ring \(R\) and a submodule \(p \subseteq M\), the quotient
  module \(M/p\) coincides, as an abelian group, with the quotient of \((M, +)\) by the
  underlying additive subgroup of \(p\), compatibly with the two quotient maps.
\end{lemma}
\begin{proof}
  \leanok
  The two quotients are defined by the same equivalence relation
  (\(x \sim y \iff x - y \in p\)), so the identity map is the required isomorphism.
\end{proof}

\begin{lemma}[Restriction intertwines the structure constants]
  \label{pa-lem:resHom_overAlgebraMap}
  \lean{AlgebraicGeometry.TwoCover.resHom_overAlgebraMap_left,
    AlgebraicGeometry.TwoCover.resHom_overAlgebraMap_right}
  \leanok
  \dcref{custom}

  \uses{pa-def:overAlgebraMap}
  For a scheme \(X\) over \(\Spec k\), opens \(U_0, U_1\) and \(a \in k\), restricting
  the structure constant of a chart to the overlap gives the structure constant of the
  overlap: \(\rho_i\bigl(a \cdot 1_{\Gamma(U_i)}\bigr) = a \cdot
  1_{\Gamma(U_0 \cap U_1)}\) for \(i = 0, 1\), where \(a \cdot 1\) denotes the image of
  \(a\) under the structure map of \ref{pa-def:overAlgebraMap}.
\end{lemma}
\begin{proof}
  \leanok
  The structure map of \ref{pa-def:overAlgebraMap} is compatible with restriction: it is
  defined by pulling back the global section \(a\) of the base, and pullback commutes
  with restriction to a smaller open.
\end{proof}

\begin{definition}[The two-cover units reduction]
  \label{pa-def:twoCover_unitsReduction}
  \lean{AlgebraicGeometry.TwoCover.unitsReduction,
    AlgebraicGeometry.TwoCover.unitsReduction_def}
  \leanok
  \dcref{custom}

  \uses{pa-def:cech_units_reduction, pa-def:moduleKSheaf}
  The \v{C}ech units reduction \(\pi\) of \ref{pa-def:cech_units_reduction}, instantiated
  at the two-chart datum of the cover: \(B = \Gamma(X, U_0 \cap U_1)\),
  \(\rho_i\) the restrictions from \(\Gamma(X, U_i)\).
\end{definition}

\begin{definition}[Truncated-exponential kernel classes of overlap sections]
  \label{pa-def:twoCover_truncExpClass}
  \lean{AlgebraicGeometry.TwoCover.truncExpClass, AlgebraicGeometry.TwoCover.kerVal,
    AlgebraicGeometry.TwoCover.kerVal_eq_coe,
    AlgebraicGeometry.TwoCover.truncExpClass_val,
    AlgebraicGeometry.TwoCover.truncExpClass_eq_engine}
  \leanok
  \dcref{custom}

  \uses{pa-def:twoCover_unitsReduction, pa-lem:cech_kernel_truncExp_class}
  For an overlap section \(b \in \Gamma(X, U_0 \cap U_1)\), the class
  \(\mathrm{tCl}(b) := [\mathrm{tExp}(b)] \in \ker \pi\)
  (\ref{pa-lem:cech_kernel_truncExp_class}), regarded in the additive notation of the
  kernel; \(\mathrm{kerVal}\) denotes the underlying thickened \v{C}ech class of an
  element of \(\ker \pi\), so \(\mathrm{kerVal}(\mathrm{tCl}(b)) =
  [\mathrm{tExp}(b)]\). Under the identification of
  \ref{pa-thm:truncExp_cech_kernel_equiv}, \(\mathrm{tCl}(b)\) is the image of the residue
  class of \(b\).
\end{definition}

\begin{lemma}[Properties of the kernel classes]
  \label{pa-lem:truncExpClass_props}
  \lean{AlgebraicGeometry.TwoCover.truncExpClass_add,
    AlgebraicGeometry.TwoCover.truncExpClass_eq_zero_iff,
    AlgebraicGeometry.TwoCover.truncExpClass_surjective}
  \leanok
  \dcref{custom}

  \uses{pa-def:twoCover_truncExpClass, pa-thm:truncExp_cech_kernel_equiv}
  \(\mathrm{tCl}(b + c) = \mathrm{tCl}(b) + \mathrm{tCl}(c)\);
  \(\mathrm{tCl}(b) = 0\) iff \(b\) is an additive coboundary
  \(\rho_1(a_1) + \rho_2(a_2)\) of the two charts; and every element of \(\ker \pi\) is
  \(\mathrm{tCl}(b)\) for some overlap section \(b\).
\end{lemma}
\begin{proof}
  \leanok
  All three restate \ref{pa-thm:truncExp_cech_kernel_equiv} through the identification of
  \ref{pa-def:twoCover_truncExpClass}: additivity and the vanishing criterion are the
  well-definedness and injectivity of the induced map, and surjectivity is its
  surjectivity.
\end{proof}

\begin{lemma}[The range of the difference map is the additive coboundary subgroup]
  \label{pa-lem:range_diff_coboundaries}
  \lean{AlgebraicGeometry.TwoCover.range_diff_toAddSubgroup,
    AlgebraicGeometry.TwoCover.h1CokAddEquivCechQuot,
    AlgebraicGeometry.TwoCover.h1CokAddEquivCechQuot_mk}
  \leanok
  \dcref{custom}

  \uses{pa-def:TwoCover_diff, pa-def:TwoCover_H1Cok, pa-def:cech_coboundary_add,
    pa-lem:submodule_quotient_addEquiv}
  The range of the restriction-difference map \(\mathrm{diff}\) of
  \ref{pa-def:TwoCover_diff}, viewed as an additive subgroup of
  \(\Gamma(X, U_0 \cap U_1)\), is exactly the additive coboundary subgroup
  \(\mathcal B^+(\rho_1, \rho_2)\) of the restriction datum
  (\ref{pa-def:cech_coboundary_add}). Consequently the two-cover cokernel
  \(\mathrm{H1Cok}\) (\ref{pa-def:TwoCover_H1Cok}) is, as an abelian group, the \v{C}ech
  quotient \(\Gamma(X, U_0 \cap U_1) / \mathcal B^+(\rho_1, \rho_2)\), compatibly with
  the quotient maps.
\end{lemma}
\begin{proof}
  \leanok
  A difference \(\rho_1(a_1) - \rho_2(a_2)\) and a sum \(\rho_1(a_1) + \rho_2(a_2)\)
  span the same additive subgroup (\ref{pa-def:cech_coboundary_add}), so the two subgroups
  coincide elementwise. The quotient statement combines this with
  \ref{pa-lem:submodule_quotient_addEquiv} applied to the submodule
  \(\mathrm{range}(\mathrm{diff})\).
\end{proof}

\begin{theorem}[\(H^1\) is the \(\varepsilon\)-kernel of the two-cover]
  \label{pa-thm:h1_truncExp_kernel}
  \lean{AlgebraicGeometry.TwoCover.h1AddEquivTruncExpCechKernel}
  \leanok
  \source{kleiman-picard}

  \uses{pa-thm:TwoCover_h1CokEquiv, pa-lem:range_diff_coboundaries,
    pa-thm:truncExp_cech_kernel_equiv, pa-def:twoCover_unitsReduction, pa-def:HModule,
    pa-def:moduleKSheaf}
  \dcref{custom}
  Let \(X\) be a scheme over \(\Spec k\) and \(U_0, U_1\) affine opens with
  \(U_0 \cup U_1 = X\). Then the degree-one cohomology of the structure sheaf is,
  additively, the kernel of the dual-number \v{C}ech units reduction of the cover:
  \[
    H^1\bigl(X, \struct{X}\bigr) \;\cong\;
    \ker\Bigl( \pi :
      \check H^1{}^\times\bigl(\Gamma(U_0 \cap U_1)[\varepsilon]\bigr) \to
      \check H^1{}^\times\bigl(\Gamma(U_0 \cap U_1)\bigr) \Bigr),
  \]
  the two-cover \v{C}ech form of
  \(H^1(C, \struct C) \cong \ker(\Pic(C_\varepsilon) \to \Pic(C))\).
\end{theorem}
\begin{proof}
  \leanok
  Compose three identifications: the two-affine-cover cokernel bridge
  \(H^1(X, \struct X) \cong \Gamma(U_0 \cap U_1)/\mathrm{range}(\mathrm{diff})\)
  (\ref{pa-thm:TwoCover_h1CokEquiv}); the carrier translation
  \(\Gamma(U_0 \cap U_1)/\mathrm{range}(\mathrm{diff}) \cong
  \Gamma(U_0 \cap U_1)/\mathcal B^+(\rho_1, \rho_2)\)
  (\ref{pa-lem:range_diff_coboundaries}); and the truncated-exponential kernel computation
  \(\Gamma(U_0 \cap U_1)/\mathcal B^+(\rho_1, \rho_2) \cong \ker \pi\)
  (\ref{pa-thm:truncExp_cech_kernel_equiv}).
\end{proof}

\begin{lemma}[Generator formula]
  \label{pa-lem:h1_truncExp_kernel_delta}
  \lean{AlgebraicGeometry.TwoCover.h1AddEquivTruncExpCechKernel_delta}
  \leanok
  \dcref{custom}

  \uses{pa-thm:h1_truncExp_kernel, pa-def:TwoCover_delta, pa-lem:TwoCover_h1CokEquiv_delta,
    pa-def:twoCover_truncExpClass}
  The identification of \ref{pa-thm:h1_truncExp_kernel} sends the connecting class
  \(\mathrm{delta}(s)\) of an overlap section \(s\) (\ref{pa-def:TwoCover_delta}) to the
  truncated-exponential class \(\mathrm{tCl}(s)\).
\end{lemma}
\begin{proof}
  \leanok
  The cokernel bridge sends \(\mathrm{delta}(s)\) to the residue class of \(s\)
  (\ref{pa-lem:TwoCover_h1CokEquiv_delta}); the carrier translation preserves residue
  classes; and the kernel computation sends the residue class of \(s\) to
  \(\mathrm{tCl}(s)\) (\ref{pa-def:twoCover_truncExpClass}).
\end{proof}

\begin{lemma}[Scaled generators]
  \label{pa-lem:smul_delta}
  \lean{AlgebraicGeometry.TwoCover.smul_delta,
    AlgebraicGeometry.TwoCover.h1AddEquivTruncExpCechKernel_smul_delta}
  \leanok
  \dcref{custom}

  \uses{pa-def:TwoCover_delta, pa-def:overAlgebraMap, pa-lem:h1_truncExp_kernel_delta}
  For \(a \in k\) and an overlap section \(s\), the \(k\)-action on \(H^1\) hits the
  connecting classes as multiplication by the structure constant:
  \(a \cdot \mathrm{delta}(s) = \mathrm{delta}\bigl((a \cdot 1) s\bigr)\); consequently
  the identification of \ref{pa-thm:h1_truncExp_kernel} sends
  \(a \cdot \mathrm{delta}(s)\) to \(\mathrm{tCl}\bigl((a \cdot 1) s\bigr)\).
\end{lemma}
\begin{proof}
  \leanok
  The connecting map \(\mathrm{delta}\) is \(k\)-linear
  (\ref{pa-def:TwoCover_delta}), and the \(k\)-action on the overlap sections is
  multiplication by the structure constant \(a \cdot 1\)
  (\ref{pa-def:overAlgebraMap}). The second claim applies
  \ref{pa-lem:h1_truncExp_kernel_delta}.
\end{proof}

\begin{definition}[The Mumford scaling of the thickened \v{C}ech group]
  \label{pa-def:twoCover_mumfordScaling}
  \lean{AlgebraicGeometry.TwoCover.mumfordScaling,
    AlgebraicGeometry.TwoCover.mumfordScaling_mk_truncExpUnit}
  \leanok
  \dcref{custom}

  \uses{pa-lem:cech_scaling_equivariance, pa-lem:resHom_overAlgebraMap}
  For \(a \in k\), the scaling \(\mu_{a \cdot 1}\) of the overlap ring's dual numbers,
  with \(a \cdot 1\) the structure constant of the overlap, descends to an endomorphism
  of the thickened \v{C}ech group of the two-cover: the structure constant is
  chart-compatible by \ref{pa-lem:resHom_overAlgebraMap}, so
  \ref{pa-lem:cech_scaling_equivariance} applies; on truncated-exponential classes it acts
  by \([\mathrm{tExp}(b)] \mapsto [\mathrm{tExp}((a \cdot 1)\, b)]\).
\end{definition}

\begin{theorem}[Mumford-scaling equivariance of the keystone]
  \label{pa-thm:mumfordScaling_equivariance}
  \lean{AlgebraicGeometry.TwoCover.mumfordScaling_h1AddEquivTruncExpCechKernel_delta}
  \leanok
  \source{kleiman-picard}

  \uses{pa-thm:h1_truncExp_kernel, pa-lem:smul_delta, pa-def:twoCover_mumfordScaling,
    pa-lem:h1_truncExp_kernel_delta}
  \dcref{custom}
  Through the identification of \ref{pa-thm:h1_truncExp_kernel}, the \(k\)-scalar action
  on \(H^1(X, \struct X)\) corresponds to the Mumford scaling of the thickened
  \v{C}ech group: for \(a \in k\) and an overlap section \(s\), the underlying
  thickened class of the image of \(a \cdot \mathrm{delta}(s)\) equals the Mumford
  scaling by \(a\) of the underlying thickened class of the image of
  \(\mathrm{delta}(s)\). Since the connecting classes generate \(H^1\)
  (\ref{pa-thm:TwoCover_delta_surjective}) and both actions are additive, this determines
  the correspondence on all of \(H^1\).
\end{theorem}
\begin{proof}
  \leanok
  \uses{pa-thm:TwoCover_delta_surjective}
  By \ref{pa-lem:smul_delta}, the image of \(a \cdot \mathrm{delta}(s)\) is
  \(\mathrm{tCl}((a \cdot 1) s)\), with underlying class
  \([\mathrm{tExp}((a \cdot 1) s)]\); by \ref{pa-lem:h1_truncExp_kernel_delta} and
  \ref{pa-def:twoCover_mumfordScaling}, the Mumford scaling by \(a\) of the underlying
  class \([\mathrm{tExp}(s)]\) of the image of \(\mathrm{delta}(s)\) is the same
  \([\mathrm{tExp}((a \cdot 1) s)]\).
\end{proof}
